\documentclass[12pt,a4paper]{article}
\usepackage[utf8]{inputenc}
\usepackage{lmodern}
\usepackage[T1]{fontenc}
\usepackage[finnish]{babel}

\usepackage{graphicx}
\usepackage{epstopdf}

\begin{document}
\title{Tietokantojen perusteet, kevät 2020 \\
Harjoitustyön raportti}
\author{Ville Manninen \\014922187 \\Macville}
\date{\today}
\clearpage\maketitle
\begin{center}
\vspace{2cm}
\includegraphics{university-of-helsinki-2.eps} 
\end{center}
\thispagestyle{empty}

\newpage
\tableofcontents
\newpage
\section{Esittely}
\subsection*{Yleistä}
Harjoitustyö on toteutettu Java-ohjelmana käyttäen NetBeans ympäristöä, ohjelmalla on seuraavat ominaisuudet.
\begin{itemize}
\item Tietokannan ja puuttuvien taulukoiden luominen.
\item Asiakkaiden lisääminen tietokantaan nimellä.
\item Paikan lisäämien tietokantaan nimellä.
\item Paketin lisääminen tietokantaan käyttäen asiakkaan nimeä.
\item Tapahtuman lisääminen tietokantaan käyttäen paikannimeä ja paketin seurantakoodia.
\item Paketin tapahtumien hakeminen
\item Asiakkaan pakettien hakeminen ja yksittäisten pakettien tapahtumien lukumäärä
\end{itemize}
\subsection*{Komennot}
Ohjelmalle annetaan komentoja komentoriviltä käyttäen pelkästään komennon numeroa, komennon jälkeen ohjelma kysyy käyttäjältä tarvittavia jatkosyötteitä. Ohjelmalla on seuraavat komennot jotka esiintyvät myös ohjelman valikossa.
\begin{description}

\item \textbf{1} Uuden asiakkaan lisääminen
\item \textbf{2} Uuden paikan lisääminen
\item \textbf{3} Paketin lisääminen
\item \textbf{4} Uuden tapahtuman lisääminen
\item \textbf{5} Tulostaa halutun paketin tapahtumat
\item \textbf{6} Tulostaa halutun asiakkaan paketit ja niihin liittyvien tapahtumien määrän
\item \textbf{0} Sulkee ohjelman 

\end{description}



\newpage
\section{Kaaviot ja Skeemat}
\subsection*{Tietokantakaavio}


\subsection*{SQL-skeema}

\newpage
\section{Tehokkuustesti}
\subsection*{Ilman indeksejä}
\subsection*{Indeksien lisäämisen jälkeen}

\newpage
\section{Tietokannan eheys}
Tietokantaan ei ole mahdollista luoda samannimisiä asiakkaita tai paikkoja, eikä tietokantaan luoduilla paketeilla saa antaa samaa seurantakoodia. Samannimisiä syötteitä valvotaan SQliten toiminnolla UNIQUE, joka varmistaa että taulukkoon ei lisätä samannimisiä arvoja. 

Ohjelma tarkistaa käyttäjän antamat syötteet ja ilmoittaa syyn jos syöte on puutteellinen, se tarkistaa myös että sille annetut syötteet eivät ole tyhjiä eikä tietokantaan lisätä NULL arvoja. Tietokantaan lisättäessä paketteja tai tapahtumia jotka viittaavat toisiin taulukoihin, sovellus varmistaa että annetut viitteet löytyvät tietokannasta.





\newpage
\section{Lähdekoodi}
\subsection*{GitHub}
\textbf{https://github.com/Viltska/sql-app-training}
\subsection*{Koodi}


\newpage

\end{document}